\documentclass[10pt,twocolumn]{article}

\usepackage[a4paper,top=17mm,bottom=19mm,left=16mm,right=16mm,columnsep=7mm]{geometry}
\usepackage[T1]{fontenc}
\usepackage{newtxtext,newtxmath}
\usepackage{microtype}
\usepackage{amsmath}
\usepackage{hyperref}
\usepackage{xcolor}

\hypersetup{colorlinks=true,allcolors=blue!55!black}
\setlength{\parindent}{1em}
\setlength{\parskip}{0pt}

\title{Cross-Module Integration and Pipeline Contracts}
\author{Liu Bowen}
\date{}

\begin{document}
\maketitle

% The final paper has one shared abstract, prepared by the team. This personal
% contribution therefore does not include a separate abstract.

\setcounter{section}{2}
\section{Discussion}

\subsection{Cross-module limitations}
\label{sec:cross-module-limitations}

The integration contract improves provenance but does not remove biological or
technical limitations. The module scores have different estimands and native
scales, so their co-location in one table is not numerical harmonization. The
number of independent GSE176206 animals is small, cell counts are unequal, and
unknown-animal observations are descriptive sensitivity data. Risk interpretation
also depends on the released gene lists, coverage rules, reference thresholds,
aggregation policy, and model version. Finally, Geneformer embedding shifts
cannot yet be used as direct inputs to the expression-based scorers.

\subsection{Future work: a multi-dataset integration framework}
\label{sec:future-multidataset-integration}

Future work should extend the contract to additional datasets and animal-condition
structures through versioned intake schemas, explicit biological-unit keys, and
machine-checkable provenance and quality-control fields. Cross-assay
transportability should remain distinct from independent external validation.
For the prediction arm, a Geneformer-derived state should be connected to the
scoring axes only after a validated predicted-expression interface is available.
This would permit in-silico screening of perturbation conditions followed by
multi-axis state evaluation, while longitudinal functional and safety endpoints
would still be required before a universal safe region could be claimed.

\setcounter{section}{3}
\section{Methods}

\subsection{Data architecture and pipeline contracts}
\label{sec:data-architecture-contracts}

The integration layer was designed to preserve the native analysis contract of
each module. Cells were treated as measurement units, whereas donors or animals
were treated as the biological inference units. This distinction was required
because arithmetic aggregation of cell-level scores and scoring of aggregated
expression profiles are different estimands. Accordingly, the integration layer
did not refit any Youth, Identity, or Risk model and did not transform their raw
score scales.

For each module, the contract recorded the input representation, preprocessing
steps, scoring unit, aggregation rule, model version, quality-control fields, and
source artifact. The cell-level Youth model retained its native scoring and donor
summary. The v3.1 donor-pseudobulk Youth framework retained M1 as the primary
operational model and M4 as a separate high-stringency sensitivity model. The two
operating points were never averaged or ensembled. The Identity and Risk modules
were likewise retained as separate named outputs.

The canonical biological-unit identifier comprised four fields:
\begin{center}
\small\ttfamily
dataset\_id \textbar{} age\_group \textbar{}\\
exact\_treatment\_arm \textbar{} animal\_label
\end{center}
\noindent Exact treatment-arm labels were retained throughout the pipeline.
Repeated numeric animal labels were not interpreted as evidence of pairing,
because animal labels were nested within treatment arms in the GSE176206
application. A pairing status was recorded only when supported by an explicit
sample map; otherwise the comparison was labelled unpaired or unknown.

The canonical output schema required the biological-unit identifier, dataset and
condition fields, module name and version, native score name, aggregation rule,
coverage and quality-control fields, evidence level, and source file. Missing
module outputs were left missing and were not imputed. Unknown-animal rows were
retained only as descriptive sensitivity observations and were excluded from the
primary known-animal inference layer.

The Risk definition was represented by two independent axes, namely
\textit{Pluripotency Risk} and \textit{Genomic Stress Risk}. Inflammation was not
included in the formal final output. The versioned gene list, scoring scripts,
and donor-summary table were registered as source artifacts; inclusion of Risk
values in the primary integrated table additionally required the input-data
reference, software environment, coverage and threshold rules, unknown-animal
policy, and clean-checkout reproducibility record.

Geneformer was represented as a separate perturbation/state-transition prediction
arm. Its embedding shift is not a predicted expression matrix and therefore
cannot be passed directly to an expression-based Youth, Identity, or Risk
scorer. Any future connection would require an independently validated
predicted-expression adapter.

\end{document}
